%%%%%%%%%%%%%%%%%%%%%%%%%%%%%%%%%%%%%%%%%%%%%%%%%%%%%%%%%%%%
%% CHAPTER 12 -- How You Talk to a Model Today
%%%%%%%%%%%%%%%%%%%%%%%%%%%%%%%%%%%%%%%%%%%%%%%%%%%%%%%%%%%%

\chapter{How You Talk to a Model Today}
\label{chap:agents}
\label{ntg:ch12}

\scenelabel
\begin{scenestrip}
\sceneitem{The problem}
\sceneitem{How we got here}
\sceneitem{Where we stand}
\sceneitem{What goes wrong}
\end{scenestrip}

\paragraph{The problem.} The previous eleven chapters described how to build the
artefact at the centre of a Large Language Model product. This
chapter describes how engineers actually use one. The two are
not the same.

In 2026, a developer who wants to build a product on top of a
frontier model does not download the model's parameters and
run them on their own hardware. They send a request, formatted
as a small piece of structured text, to a server operated by
the model's vendor. They receive a response, sometimes streamed
back word by word as the model produces it. They may, in the
course of a single user interaction, let the model invoke
external tools that they have described to it in a structured
schema. They may route the whole interaction through a
standardised protocol that lets the model use tools written by
third parties. The result is wrapped in a software development
kit (SDK) and embedded inside the application the user actually
sees.

Each of those layers is the result of a deliberate engineering
decision. Each layer has its own failure modes. Each layer is
the place where one team's bug becomes another team's incident
report. The problem of this chapter is the surface between the
model and the application that uses it, and the discipline
required to operate that surface in production without
incident. The non-technical reader does not need to understand
the SDKs in detail. The non-technical reader does need to know
that the surface exists, what kinds of things move across it,
and what kinds of things go wrong when they move incorrectly.

\paragraph{How we got here.} The first widely available LLM API was OpenAI's GPT-3
completions endpoint in 2020. The contract was simple. A
developer sent a string of text. The endpoint returned a
string of text. There was no notion of conversation, no
notion of tool use, no notion of structured output. The model
was a sophisticated autocomplete, and the API treated it as
such.

Within two years the contract had grown. The chat-completions
schema, introduced in early 2023, distinguished between
\emph{system} messages (which set the model's role and
instructions), \emph{user} messages (which carried the user's
input), and \emph{assistant} messages (which carried the
model's responses). This made multi-turn conversations natural
to express. The function-calling extension, introduced in
mid-2023, let the model emit structured JSON conforming to a
developer-supplied schema, which the surrounding application
could validate and dispatch as a function call. Streaming over
server-sent events let the model's response arrive token by
token rather than as a single block, so users could see
responses being typed in real time. Prompt caching let
identical prefixes of repeated requests be re-used without
recomputation, reducing both cost and latency.

By 2024 most providers --- OpenAI, Anthropic, Google, xAI,
Cohere, Mistral, the open-source serving ecosystem --- exposed
an OpenAI-compatible endpoint. The interface had won by
default. A developer who learned to call OpenAI's API could
call any of the others with a few lines of configuration
change.

The next layer was \emph{agents}. The earliest agentic systems
--- ReAct from Princeton and Google in early 2023, AutoGPT and
BabyAGI from independent developers later that year --- were
demonstrations in which a single control loop drove a language
model to take actions in service of a goal. The model was
asked, repeatedly: ``you have these tools available, you have
this goal, what would you like to do next?'' The model
responded with a tool call. The tool call was executed. The
result was returned to the model. The model decided whether it
was done or whether it wanted to take another action. The
loop continued until the model reported completion or the
operator stopped it.

The frameworks that followed --- LangChain, LangGraph, AutoGen,
CrewAI, smolagents, the Anthropic Agent SDK --- codified the
pattern. They added memory, retry logic, branching, multi-
agent collaboration, and the operational machinery to make
agent loops survive contact with production. In parallel,
IDE-integrated agents --- GitHub Copilot, Cursor, Claude Code
--- moved from autocomplete demos to systems that could read
an entire codebase, design a change spanning many files,
write the change, run the tests, and commit the result.

The most recent layer is the \emph{Model Context Protocol},
proposed by Anthropic in late 2024 and adopted across the
ecosystem in 2025. MCP is a standardised protocol for how a
host application (a chat client, an IDE, an SDK harness)
communicates with a tool server. Before MCP, every host had
to write a custom integration for every tool, and the
combinatorial cost was substantial. After MCP, every host
speaks the same protocol, every tool implements the same
protocol, and integration becomes plumbing. Standardisation
mattered, in the same way that standardisation has always
mattered in computing.

\paragraph{Where we stand.} A modern LLM application is a layered stack. At the bottom are
the model parameters and the inference runtime that runs them.
Above that is the wire protocol, typically OpenAI-compatible
chat completions with streaming. Above that is the SDK --- a
typed client library, in the developer's preferred programming
language, that handles retries, rate limiting, structured-
output validation, and cost accounting. Above that is the
orchestration layer --- the agent framework, the retrieval
pipeline, the memory store, the session management. At the
top is the host application, which is the only thing the user
ever sees. The Model Context Protocol cuts across the
orchestration layer, providing a standard way for the host to
discover and invoke tools written by third parties.

The whole stack is designed for change. Any layer can be
replaced, in principle, without disturbing the others. A team
can switch model providers by changing a configuration value.
A team can add a tool by deploying an MCP server. A team can
swap out the orchestration framework by rewriting the layer
above the SDK. None of this is free in practice, but the
modularity is real, and it is the property that has let the
ecosystem evolve as fast as it has.

\paragraph{What goes wrong.} The transition from research demonstration to production
system is where the failure modes accumulate.

Rate limits turn into cascading retries that overwhelm the
provider; this is the classical \emph{thundering herd}
problem from distributed systems, with a new wrapper. Tool
schemas drift between development and production, so a tool
the model has been trained to call no longer accepts the JSON
the model emits. Streaming responses arrive out of order, or
fail mid-stream, or are silently truncated by an intermediate
proxy. Prompt caches return stale results when the cache key
representation changes between client versions.

Agent loops can enter pathological cycles, repeating the same
tool call until the budget is exhausted. The cost of an
unbounded agent run can reach four figures before anyone
notices. There are stories, increasingly less apocryphal, of
agents that ran overnight and produced bills that exceeded
their entire month's budget. The fix is a hard limit on
iterations, on tokens, on tool calls, and on dollars. Not
every team adds the limits before they need them.

The MCP standard helps with integration but introduces new
trust boundaries. A malicious MCP server can exfiltrate data,
lie about tool results, or inject instructions into the
model's context. The discipline of vetting MCP servers is, at
this writing, immature, and the analogy to the early days of
browser plug-ins is uncomfortably apt. The \emph{prompt
injection} attack vector grows alongside every new capability
the agent gains, and a system that can read documents, write
files, and call APIs has many more surfaces on which an
injection can land than a chat assistant does.

\section{The idea, in plain language}

\subsection*{The agent loop}

The clearest way to understand a modern agentic LLM system is
to look at the loop that drives it. The loop is conceptually
simple.

\begin{enumerate}
  \item The user gives the system a goal.
  \item The system constructs a prompt that contains the goal,
        the available tools (each described by a schema), and
        any prior context.
  \item The system asks the model what it would like to do
        next. The model returns either a final answer to the
        user or a tool call.
  \item If the model returned a tool call, the system
        validates the call against the tool's schema, executes
        the tool, and appends the result to the prompt.
  \item Return to step 3.
  \item When the model returns a final answer, return that
        answer to the user.
\end{enumerate}

That is the entire loop. It is a few hundred lines of code in
a typical implementation. The complexity of an agentic
system is not in the loop itself but in the surrounding
infrastructure: the tool schemas, the memory store, the cost
controls, the error handling, the safety filters, the
observability, the user interface.

The loop has several characteristic operating modes. In the
\emph{single-shot} mode, the model returns a final answer
immediately and the loop terminates. This is what happens when
a user asks a chat assistant a question that does not require
any tools. In the \emph{tool-call} mode, the model takes one or
more tool actions before returning a final answer. This is
what happens when an assistant looks up a customer record, or
runs a calculation, or sends an email. In the
\emph{multi-step planning} mode, the model returns a sequence
of intended actions before executing any of them, and the
system asks for confirmation. This is what happens in some
coding agents, which propose a multi-file change for review
before applying it.

The boundaries between these modes are not crisp, and a single
deployed system may move between them depending on the task.

\subsection*{The Model Context Protocol, briefly}

The Model Context Protocol is, structurally, a set of
conventions for how three parties --- a host application, a
client running inside the host, and a tool server running
elsewhere --- talk to one another. The conventions cover the
discovery of available tools (``what can you do?''), the
invocation of tools (``please do this''), the return of
results (``here is what happened''), the management of
sessions, and the negotiation of capabilities and security.

The appeal of the protocol is that it dissolves what
engineers call the \emph{M times N integration problem}.
Without a standard, every host application that wanted to use
every tool would have to write a custom integration for it,
and the number of integrations would grow as the product of
the number of hosts and the number of tools. With a standard,
every host speaks the standard, every tool implements the
standard, and integrations become plumbing rather than custom
engineering.

The non-technical implication is that the ecosystem of
LLM-accessible tools is, as of 2025, growing rapidly. A user
of Claude Desktop, Cursor, or any other MCP-aware host can
add capabilities to the assistant simply by enabling the
corresponding MCP server: a connector for the user's
calendar, a connector for their database, a connector for
their company's internal documentation. The convenience is
real. So is the security exposure.

\subsection*{Trust boundaries in agentic systems}

The single most important conceptual idea for a non-technical
reader to take away from this chapter is the notion of a
\emph{trust boundary}.

A trust boundary is a place in the system where information
crosses from a less-trusted source to a more-trusted one, and
where verification is required to maintain the integrity of
the more-trusted side. In an agentic LLM system, the trust
boundaries are numerous.

The user's input is, in general, less trusted than the
system's instructions. The model has to be told not to follow
user-supplied instructions that contradict its system prompt,
and the system has to be designed so that user-supplied input
cannot be confused with system-supplied input.

The output of a tool is, in general, less trusted than the
inputs the model gave to the tool. A tool that fetches a web
page may return content that contains an instruction the
model should not follow. A tool that queries a database may
return data that contains a malicious payload. The system has
to treat tool outputs as data rather than as instructions.

The output of one agent in a multi-agent system is, in
general, less trusted by the other agents than the original
user's request. A system in which agents pass instructions to
one another can be subverted at any link in the chain.

Maintaining the trust boundaries is, in many production
deployments, the single largest piece of operational work.
The model itself does not enforce them; the surrounding
system has to. A system that does not enforce them is a
system that can be manipulated through any of the inputs it
accepts, which is to say, every system that has not been
carefully designed.

\section{The engineering consequence}

\paragraph{Costs are real and need bounding.} The per-query
cost of a single chat-completion call is small. The cost of
an agentic system that can call itself recursively is, in
principle, unbounded. Production deployments need hard limits
on iterations, on tokens consumed, on tool calls made, and on
dollars spent. Without those limits, the system's worst
behaviour is the bill at the end of the month. With the
limits, the system has well-defined failure modes that can be
detected and recovered from.

\paragraph{Latency budgets compose.} Each layer of the stack
adds latency: the network round-trip, the model's
inference time, the tool's execution time, the orchestration
overhead. An agent loop that calls three tools sequentially
incurs the latency of each. A user who has clicked a button
and is waiting will not, in general, wait many seconds.
Designing an agentic system to meet a user-experience
latency budget requires conscious decisions about which tools
to call in parallel, which to call asynchronously, and which
to skip when the budget is tight.

\paragraph{Vendor lock-in is real but manageable.} An
application built tightly against one vendor's specific API
extensions will be expensive to port to another vendor.
Applications that use the OpenAI-compatible interface and
avoid vendor-specific features can usually be moved to
another vendor with a configuration change. The trade-off is
between using each vendor's most powerful features and
preserving the option to switch. Different teams sit in
different places on the trade-off. The discipline is to make
the choice consciously rather than by accident.

\paragraph{Verification becomes the binding constraint as
generation cheapens.} When language models were slow and
expensive, generation was the bottleneck in any LLM workflow.
As inference has become faster and cheaper, that bottleneck
has shifted. A team using an LLM to draft contracts, code, or
analysis at high volume will find that the rate-limiting step
is no longer generating content but verifying it. This is not
a failure mode. It is a structural shift in where the work
lives, and it has a straightforward implication: invest in
evaluation and verification infrastructure early, before the
volume of generated content exceeds the team's capacity to
review it. A product that generates faster than it can verify
will eventually ship errors that a slower product would have
caught.

\paragraph{Observability is non-negotiable.} A production
LLM system needs to log every request, every response, every
tool call, and every error. The logs need to be searchable,
auditable, and retained for as long as regulatory and
operational needs require. The OpenTelemetry community has
developed conventions for instrumenting LLM applications;
they are increasingly the de facto standard. A team that
deploys an LLM application without observability is a team
that will not understand its own system's behaviour when
something goes wrong.

\section{What goes wrong, and why}

\paragraph{Prompt injection.} The single most consequential
class of attack on agentic LLM systems. An attacker crafts an
input --- a document, a webpage, an email, a tool result ---
that contains an instruction the model is induced to follow.
The instruction may be to exfiltrate data, to perform an
action the user did not authorise, or to ignore the system
prompt. The mitigations are layered: input sanitisation,
prompt structure that flags untrusted content, output
filtering, monitoring for anomalous behaviour. None of the
mitigations is complete. As of this writing, prompt injection
is the unsolved problem of agentic deployment, and any team
building a production agent should treat it as a category of
risk that requires defence in depth, not a single fix.

\paragraph{Confused-deputy attacks.} A particularly insidious
form of prompt injection in which an attacker uses the
model's privileges to do something the attacker themselves
could not. The model has access to internal systems; the
attacker tricks the model into using that access on the
attacker's behalf. The classic computer-security pattern of
the confused deputy --- the deputy has the authority, but is
acting on instructions from someone who does not --- now has
an LLM-shaped instance, and the defences are similarly
familiar: principle of least privilege, mediation of
sensitive actions, and audit.

\paragraph{Cost runaways.} An agent loop with no termination
condition, or with a poorly-tuned one, can iterate indefinitely
and produce a bill that exceeds the operator's budget by
orders of magnitude. The fix is hard limits: a maximum number
of iterations, a maximum number of tool calls, a maximum
spend per session, and an alerting system that notifies an
operator before any of those limits is approached. The fix is
mundane. It is also reliably effective when actually
implemented, and reliably absent in the early versions of
many agentic systems.

\paragraph{Agent loops that go in circles.} A model that has
been asked to accomplish a task, encountered an obstacle, and
not figured out how to get around the obstacle will sometimes
repeat the same set of tool calls indefinitely, hoping for a
different result. The mitigations include better prompts that
encourage the model to recognise repeated failure, hard
limits on retries of the same call, and monitoring that
detects loops and intervenes. None of these is fully
satisfying. The honest position is that agent loops are an
active research area, and the failure modes are still being
catalogued.

\paragraph{Vendor-side updates that break the integration.}
A vendor that updates the underlying model can change the
model's behaviour in ways that break carefully tuned tool
schemas, prompt templates, and agent loops. The deployer has
to test before any vendor-side update, has to maintain the
ability to roll back, and has to maintain a test suite that
catches the relevant kinds of regression. This is engineering
discipline applied to a moving target, and it is one of the
ongoing operational costs of building on a vendor-served
model.

\paragraph{Trust-boundary leakage in multi-agent systems.}
Systems in which several agents pass tasks and information
between one another are particularly susceptible to a class
of failure in which an instruction injected at one point in
the chain propagates to others without being filtered. The
discipline of treating each agent's output as untrusted
input to the next one is unfamiliar and, in our experience,
reliably skipped on the first design pass. The fix is to
treat the multi-agent architecture as a distributed system in
which every link is a trust boundary.

\section*{Bridges}
\addcontentsline{toc}{section}{Bridges}

This chapter built on Chapter~\ref{chap:alignment} (the aligned model that the
APIs serve), Chapter~\ref{chap:systems} (the retrieval and tool-use systems
that surround the model), and Chapter~\ref{chap:governance} (the governance
discipline within which the system operates). It described
the developer-facing surface of the technology: the layer at
which most LLM products are actually built, and the layer at
which most production failures occur. Chapter~\ref{chap:error-model} returns to
the failure modes catalogued here as the \emph{environment}
axis of the error model, where the model is at the mercy of
the prompt and the tools the surrounding system has provided.
The next, and final, chapter --- \emph{A Map of How LLMs Go
Wrong} --- gathers all the failure modes from across the
preceding twelve chapters into a single diagnostic framework
the non-technical reader can use to make sense of any
particular wrong answer.
